\subsection{相关研究与本题的建模切入点}
\label{sec:research_position}
\subsubsection{无人机配送中的路线与资源协同}
无人机配送的优化模型首先取决于运输组织方式。Otto等从运筹决策角度梳理了民用无人机的路径、部署与资源管理问题\cite{otto2018}；Macrina等进一步总结无人机辅助配送的路线模型和求解方法\cite{macrina2020}。这些研究说明，同样以无人机承担交付，固定基地直达、多个服务点连续配送和车机协同具有不同的资源连接关系，建模时应先辨明对象，再选择变量与约束。

Murray与Chu的车机协同配送模型\cite{murray2015fstsp}以及Agatz等的旅行商—无人机优化\cite{agatz2018}为同步组织提供了典型研究框架。本题的起降基地固定，不涉及移动卡车平台，故直接保留三类运输机的容量、速度和能源差异，以完整运输任务连接箱组、路线与执行时刻。近期混合无人机医疗配送研究同时考虑机队构成、额外电池与严格交付期限\cite{bhuiyan2026}，与本题的异构资源选择具有相近的决策特征。

灾后运输的评价还应与保障对象一致。Al Theeb与Murray将灾后路径与资源配送联立考虑\cite{altheeb2017relief}；Yucesoy等将无人机灾害响应研究划分为信息收集、物资交付与通信恢复等任务\cite{yucesoy2025}。本题同时含物资交付与通信保障，因此本文把箱级需求、任务能源和全过程链路统一到可执行日程，而非仅以总距离衡量救援质量。

\subsubsection{载荷能耗与电池恢复的联合决策}
Dorling等把飞行能耗与无人机路径优化联系起来\cite{dorling2017}；Zhang等比较了配送无人机能耗模型，并发布相应勘误\cite{zhang2021energy,zhang2023corrigendum}。在本题中，每站卸货改变后续航段载荷，地形爬升又构成附加成本，因而应按给定物理函数逐段计算，而不能对所有飞行使用一个平均能耗系数。

有限电池会进一步改变路线的系统代价。Liu等把电池安装、充电安排与车机配送路线联立研究\cite{liu2025battery}，强调可恢复电池库存与路线决策的相互作用。这些研究促使本文把能源恢复直接落实为任务资源约束，而不是在路线确定后附加一个统一等待时间。本文的充电位置固定，但同型电池数量有限。任务开始须同时具备兼容飞机和满电电池，返航后两类资源又在不同时刻释放。因此，物理上更节能的任务不仅降低电量消耗，也可能通过缩短恢复时间改善后续接续。

\subsubsection{连续通信与中继服务窗口}
无人机通信研究强调部署、轨迹与链路之间的耦合\cite{zeng2016uavcom,zeng2019sky}。Zhang等将蜂窝连通性纳入无人机轨迹优化\cite{zhang2019cellular}，Zeng等分析旋翼无人机通信的能源与时间配置\cite{zeng2019energy}。本题的中继除了满足接入与回传条件，还必须在运输任务经过相应位置时实际在线；几何覆盖和时间覆盖缺一不可。

Hong等在连续取送任务中施加全过程通信要求\cite{hong2025}；Akpunar与Akpinar围绕灾害中的不间断网络覆盖构建优化模型\cite{akpunar2025}。本文采用真实地形、分阶段航迹和中继在线窗口，把直连失效表示为连续需求区间，再选择服务提供方和资源顺序。连续区间及切换端点分别核验，使网络保障与运输执行保持一致。

\subsubsection{名义性能、保障水平与执行恢复}
稳健优化关注名义性能和不确定性保障之间的代价\cite{bertsimas2004}。Xin等以两阶段分布鲁棒方法研究救灾车机网络\cite{xin2025}；Dukkanci同时考虑需求和无人机旅行时间不确定下的选址—路径决策\cite{dukkanci2026}。这类研究依赖相应不确定集合或分布信息。本题没有给出实测偏差分布，本文采用确定性的压力情景，分别检验原方案、规定范围内的时刻修复及结构调整。

Li等在飓风环境下进行动态车机运输规划\cite{li2025hurricane}，Zhang等把通信中断与事件触发的滚动规划结合\cite{zhang2026communication}。它们提示，在执行信息变化时，应明确哪些决策可以调整、哪些已经固定。本文以能量、资源恢复和通信窗口为触发对象，给出相应调整代价，将敏感性分析落实为方案选择与恢复依据。

\subsection{方法组织与实质输出}
本文采用统一物理接口、分层优化和独立核验相结合的组织方式。小规模组批用精确方法闭合最少架次；大规模调度同时构造执行方案与下界；连续通信使用候选部署、整数排程和连续时间精修；固定日程配置利用区间结构完整枚举。表\ref{tab:contribution_map}概括每一层新增的决策机制和输出。
\begin{table}[htbp]\centering\small
\caption{四个决策层次的方法与实质输出}\label{tab:contribution_map}
\begin{tabularx}{\linewidth}{p{1.5cm}LLL}\toprule
层次 & 新增机制 & 计算方法 & 可核验输出\\\midrule
物理任务 & 地形、载荷、交接、返航 & 栅格相交与逐段递推 & 时间、能耗、交付偏移\\
单点组批 & 不可拆箱与机型选择 & 精确覆盖、计数状态递推 & 18架次及安全余量跃迁\\
运输调度 & 飞机—电池不同步 & 邻域搜索、事件排程、独立下界 & 完整日程及4.5541\%间隙\\
通信保障 & 连续链路与服务交接 & 候选部署、联合整数规划、时序精修 & 联合方案与保障备选\\
独立配置 & 任务不可拆、组间不共享 & 关联压缩、区间峰值、完整枚举 & 逐组配置与库存响应\\\bottomrule
\end{tabularx}
\end{table}

四问之间传递的是可核验的任务属性和选定方案，而非未经说明的实验简称。后续分析围绕“任务怎样构成、资源怎样接续、通信怎样保障、组织怎样配置”逐层展开。每一问都同时保留方案本身、质量评价和对应参数响应，从而把计算规模转化为可读的决策解释。
